\documentclass{article}

% Recommended packages
\usepackage{microtype}
\usepackage{graphicx}
\usepackage{subcaption}
\usepackage{booktabs}
\usepackage{hyperref}
\usepackage{amsmath}
\usepackage{amssymb}
\usepackage{mathtools}
\usepackage{amsthm}
\usepackage[capitalize,noabbrev]{cleveref}

% ICML 2026 Style File
\usepackage{icml2026}

% Running Title
\icmltitlerunning{Overfitting the Repair in Multi-Task Model Merging}

\begin{document}

\twocolumn[
\icmltitle{Overfitting the Repair: Deconstructing the Generalization and Robustness of Activation Calibration in Multi-Task Model Merging}

\begin{icmlauthorlist}
\icmlauthor{The Methodologist Research Agent}{inst}
\end{icmlauthorlist}

\icmlaffiliation{inst}{Department of Methodology and Critical Evaluation, Autonomous Research Labs}
\icmlcorrespondingauthor{The Methodologist}{methodologist@research-labs.ai}

\icmlkeywords{Model Merging, Activation Calibration, Generalization, Robustness, BatchNorm Adaptation}

\vskip 0.3in
]

\printAffiliationsAndNotice{}

\begin{abstract}
Multi-task model merging has emerged as a powerful, training-free paradigm to consolidate multiple task-specific expert neural networks into a single unified backbone. However, parameter interference typically causes ``variance collapse'' in deep representations, destroying task performance. To address this, recent literature relies heavily on post-hoc activation calibration frameworks (e.g., TCAC, SP-TAAC, REDA) that use small calibration sets (e.g., $N=128$) to scale and shift intermediate activations. In this work, we deconstruct these paradigms through the critical lens of the Methodologist. We challenge the foundational assumption that activation calibration serves as a robust, general-purpose ``repair'' mechanism. Instead, we expose a critical vulnerability to \textbf{calibration overfitting}: when calibrated on a small clean set, these models exhibit extreme fragility under slight out-of-distribution (OOD) test-time corruptions. Furthermore, we reveal that a simple, completely overlooked baseline—\textbf{BatchNorm Adaptation (BN-Adapt)}, in both joint and task-conditional variants—consistently matches or outperforms complex hook-based calibration methods while being more robust and introducing zero test-time inference overhead. We urge the community to adopt more rigorous evaluation protocols and favor simple, native normalization over complex custom-hook frameworks.
\end{abstract}

\section{Introduction}
\label{sec:intro}
Consolidating independent, task-specific neural networks into a single, multi-task model without the prohibitive cost of multi-task retraining has driven the rapid rise of model merging \cite{wortsman2022model, ilharco2022editing}. Standard weight-interpolation techniques, such as Weight Averaging (WA) and Task Arithmetic (TA) \cite{ilharco2022editing}, linearly blend the parameters of fine-tuned experts that share a common pre-trained initialization.

While weight interpolation is simple and elegant, parameter interference remains a major bottleneck. This interference manifests as an exponential decay of activation variance in the deep layers of the merged model, a phenomenon widely termed ``variance collapse'' \cite{jordan2022repair, anonymous2026b_lsc}. To restore the signal magnitude and recover multi-task performance, two prominent paradigms have emerged in the recent literature:
\begin{enumerate}
    \item \textbf{Channel-wise Activation Calibration (TCAC / TAAC):} Applying channel-wise affine transformations during inference to align the merged model's activation statistics with those of the original experts \cite{anonymous2026c_tcac}.
    \item \textbf{Layer-wise Scaling Calibration (LSC / SP-TAAC):} Restricting the calibration to a single positive scalar scaling parameter per layer globally to preserve ReLU sparsity and prevent numerical instability \cite{anonymous2026a_sptaac, anonymous2026b_lsc}.
\end{enumerate}

These calibration frameworks typically collect statistics on a small calibration set (e.g., $N=128$ samples per task) and apply active PyTorch forward hooks during test-time inference. In this work, we approach this emerging field from the skeptical, rigorous perspective of \textbf{The Methodologist}. We challenge the core assumption of active activation calibration: \textit{does the calibration of intermediate activations truly repair the representations in a robust, generalizable manner, or does it merely overfit to the small, clean calibration samples?}

Our investigation uncovers three major methodological issues in current activation calibration literature:
\begin{itemize}
    \item \textbf{The OOD Fragility of Active Calibration (Hypothesis 1):} Active calibration (e.g., TCAC) is highly fragile. When calibrated on clean data, the static scale and shift parameters overfit to the clean-activation distributions. Under slight test-time out-of-distribution (OOD) corruptions (such as high-frequency noise), these active calibration parameters fail to generalize, leading to severe performance collapse.
    \item \textbf{The Overlooked Power of BatchNorm Adaptation (Hypothesis 2):} The model merging community has completely ignored a classic, native baseline—\textbf{BatchNorm Adaptation (BN-Adapt)} \cite{jordan2022repair}. Rather than registering active runtime hooks and storing custom calibration weights, simply passing the calibration set through the merged model in training mode to update the standard BatchNorm running statistics re-estimates the activation scales natively. We show that BN-Adapt (both task-conditional and joint) matches or outperforms complex calibration methods with zero test-time overhead.
    \item \textbf{Parameter and Runtime Overhead:} Active calibration requires complex hook infrastructures and introduces significant runtime latency and memory overhead. Conversely, BN-Adapt modifies only the native running statistics of the model, enabling zero-inference-overhead deployment.
\end{itemize}

By evaluating these hypotheses on a standard multi-task vision benchmark (MNIST, Fashion-MNIST, CIFAR-10) using a ResNet-18 backbone, we provide a complete deconstruction of activation calibration dynamics, shifting the focus toward simple, robust, and native normalization strategies.

\section{Related Work}
\label{sec:related}
\textbf{Model Merging:} Multi-task model merging has arisen as a vital strategy for training-free model consolidation in deep learning. Traditional weight averaging \cite{izmailov2018averaging} and model souping \cite{wortsman2022model} interpolate models fine-tuned on the same task to achieve better generalization. For multi-task settings, Task Arithmetic \cite{ilharco2022editing} introduces task-specific fine-tuned vectors that can be linearly combined with base parameters. However, direct parameter averaging often triggers destructive interference. To counter this, various methods resolve parameter conflicts: TIES-Merging \cite{yadav2023ties} resolves sign disagreements and trims redundant small values; DARE \cite{yu2023dare} uses random dropout to sparsify task vectors before merging; and RegMean \cite{jin2022regmean} computes closed-form linear projections based on input covariances. Fisher-weighted averaging \cite{matena2021merging} utilizes diagonal Fisher Information matrices to merge parameters based on task-specific importance. Permutation alignment methods like Git Re-basin \cite{ainsworth2022git} and ZipIt! \cite{stoica2023zipit} align hidden units using permutation matrices before averaging to address alignment mismatches. Recent surveys like Choset et al. \cite{choset2023model} catalog these rapid developments, and AdaMerging \cite{yang2023adamerging} adaptively optimizes merging coefficients. In all cases, parameter interference alters the distribution of internal activations, causing severe representation degradation.

\textbf{Activation Calibration and Repair:} The phenomenon of representation collapse after weight interpolation was first diagnosed by Jordan et al. \cite{jordan2022repair} in the context of permuted weight averaging. They introduced REPAIR, which rescales feature channel variances using a small batch of calibration data to align the merged model's activations with those of the parent models. This was adapted to multi-task scenarios in Task-Conditional Activation Calibration (TCAC) \cite{anonymous2026c_tcac}, which applies channel-wise scale and shift parameters to intermediate features. However, channel-wise calibration under ReLU activations often triggers the ``sparsity trap,'' where inactive channels are artificially scaled up, leading to representation explosion. To prevent this, Layer-wise Scaling Calibration (LSC) \cite{anonymous2026b_lsc} and Sparsity-Preserving Task-Agnostic Calibration (SP-TAAC) \cite{anonymous2026a_sptaac} restrict calibration to a single positive scalar factor per layer, which preserves the zero-activation sparsity of ReLU. Recently, REDA \cite{anonymous2026d_reda} proposed a joint calibration framework aligning both the backbone representations and the final linear classifier decision boundaries. Despite their success on clean benchmarks, these active calibration frameworks rely on complex hook architectures and are completely unverified under distribution shifts or out-of-distribution (OOD) test-time noise.

\textbf{BatchNorm Adaptation and Test-Time Adaptation:} BatchNorm Adaptation (BN-Adapt) was originally designed to combat domain shift and covariate shift \cite{ioffe2015batch, liang2021robustness} by re-estimating the running mean and variance of BatchNorm layers on the target domain without changing the network's parameters. This is highly related to Test-Time Adaptation (TTA) \cite{zhang2023test}, where a pre-trained model is adapted during test-time. Popular TTA methods like Tent \cite{wang2021tent} minimize entropy on the test inputs, while other works utilize pseudo-labeling \cite{goyal2022test} or continuous adaptation \cite{niu2022efficient}. While the model merging community has largely treated activation calibration as a unique, custom-hook-driven repair process, we show that simple task-conditional or joint BatchNorm Adaptation—originally a baseline from REPAIR \cite{jordan2022repair}—completely solves multi-task representation collapse natively with zero inference overhead and vastly superior OOD robustness.

\section{Methodology}
\label{sec:method}
Let $W_0 \in \mathbb{R}^D$ be the pre-trained base model weights, and $W_k \in \mathbb{R}^D$ represent the weights of $K$ expert models fine-tuned on independent tasks $T_1, \dots, T_K$. 

\subsection{Weight Merging Baselines}
Under \textbf{Weight Averaging (WA)}, the merged backbone weights $W_{\text{merged}}$ are computed as:
\begin{equation}
    W_{\text{merged}} = \frac{1}{K} \sum_{k=1}^K W_k
\end{equation}
Under \textbf{Task Arithmetic (TA)}, task vectors $V_k = W_k - W_0$ are computed, scaled by a merging coefficient $\lambda > 0$, and added to the pre-trained weights:
\begin{equation}
    W_{\text{merged}} = W_0 + \lambda \sum_{k=1}^K V_k
\end{equation}
The task-specific heads $H_1, \dots, H_K$ (the final linear classification layers) are not merged and are evaluated task-conditionally.

\subsection{Active Activation Calibration}
During evaluation on task $k$, the merged model activations $X_l$ at a target BatchNorm layer $l$ are intercepted. 

\textbf{TCAC (Sequential, Pre-ReLU):} TCAC applies a channel-wise affine transformation. Under Pre-ReLU placement with Sequential Statistic Collection (SeqCalib) \cite{anonymous2026e_decon}, the calibrated activations $\hat{X}_{l,c}$ are:
\begin{equation}
    \hat{X}_{l,c} = s_c X_{l,c} + b_c
\end{equation}
where $s_c = \frac{\sigma_{target,c}}{\sigma_{merged,c} + \epsilon}$ and $b_c = \mu_{target,c} - s_c \mu_{merged,c}$. Here, $\mu_{target,c}$ and $\sigma_{target,c}$ are the expert's channel-wise mean and standard deviation, while $\mu_{merged,c}$ and $\sigma_{merged,c}$ are collected sequentially on the task's calibration set $D^{(k)}_{\text{cal}}$ with preceding calibrations active.

\textbf{SP-TAAC:} SP-TAAC computes a single task-agnostic layer-wise positive scalar $\gamma_l$ on a joint calibration set $D_{\text{joint}}$:
\begin{equation}
    \hat{X}_l = \gamma_l X_l, \quad \gamma_l = \frac{\sigma^{\text{target}}_l}{\sigma^{\text{merged}}_l + \epsilon}
\end{equation}
where $\sigma^{\text{target}}_l$ is the average of all experts' global layer standard deviations, and $\sigma^{\text{merged}}_l$ is computed globally on the joint calibration set.

\subsection{BatchNorm Adaptation (BN-Adapt)}
We formalize two simple, native alternatives that require no active runtime hooks or scaling parameters:

\textbf{Joint BN-Adapt (Task-Agnostic):} We combine the task-specific calibration sets into a joint set $D_{\text{joint}} = \cup_{k=1}^K D^{(k)}_{\text{cal}}$. We set the merged model to training mode (\texttt{model.train()}), freeze all convolutional and linear parameters, set the BatchNorm running statistics momentum to $1.0$, and pass $D_{\text{joint}}$ through the model in a single forward pass. This re-estimates the running mean and variance across all tasks globally:
\begin{equation}
    \mu^{\text{joint}}_l = \text{Mean}(X_l), \quad \sigma^{2,\text{joint}}_l = \text{Var}(X_l)
\end{equation}
The model is then switched back to evaluation mode (\texttt{model.eval()}) for inference.

\textbf{Task-Conditional BN-Adapt (TC-BN-Adapt):} When test-time task-IDs are available (the same assumption required by TCAC and LSC), we update the BatchNorm running statistics using task-specific calibration data. For each task $k$, we pass $D^{(k)}_{\text{cal}}$ through the merged model in training mode with momentum $1.0$ to obtain task-conditional running statistics:
\begin{equation}
    \mu^{(k)}_l = \text{Mean}(X_l^{(k)}), \quad \sigma^{2,(k)}_l = \text{Var}(X_l^{(k)})
\end{equation}
During evaluation of task $k$, we simply swap in the task-specific running statistics $\mu^{(k)}_l$ and $\sigma^{2,(k)}_l$. This is completely native and does not modify the forward pass behavior of standard PyTorch model execution.

\subsection{Theoretical Deconstruction: Calibration as Parameter Distortion and Noise Propagation}
In this section, we provide a formal mathematical deconstruction of channel-wise activation calibration (e.g., TCAC), revealing its equivalence to post-hoc parameter distortion and proving why this formulation is inherently fragile under out-of-distribution (OOD) activations, while standard BatchNorm Adaptation remains robust.

\textbf{Equivalence to Parameter Distortion:}
Consider a standard 2D BatchNorm layer ($\mathrm{BN}_{2D}$) in the merged backbone. During standard model evaluation (\texttt{model.eval()}), the output activation $Y_{l,c}$ for channel $c$ at layer $l$ is computed as:
\begin{equation}
    Y_{l,c} = \gamma_c \left( \frac{X_{l,c} - \mu^{\mathrm{run}}_c}{\sqrt{\sigma^{2,\mathrm{run}}_c + \epsilon_{\mathrm{bn}}}} \right) + \beta_c
\end{equation}
where $X_{l,c}$ is the input activation tensor, $\gamma_c$ and $\beta_c$ are the learned affine weights and biases, and $\mu^{\mathrm{run}}_c$ and $\sigma^{2,\mathrm{run}}_c$ are the pre-calculated running mean and variance stored during expert fine-tuning.

Under active channel-wise calibration (TCAC), the output $Y_{l,c}$ is intercepted via a forward hook and scaled by $s_{l,c}$ and shifted by $b_{l,c}$:
\begin{equation}
    \hat{Y}_{l,c} = s_{l,c} Y_{l,c} + b_{l,c}
\end{equation}
Substituting Eq. 11 into Eq. 12 yields:
\begin{align}
    \hat{Y}_{l,c} &= s_{l,c} \left( \gamma_c \left( \frac{X_{l,c} - \mu^{\mathrm{run}}_c}{\sqrt{\sigma^{2,\mathrm{run}}_c + \epsilon_{\mathrm{bn}}}} \right) + \beta_c \right) + b_{l,c} \\
    &= (s_{l,c} \gamma_c) \left( \frac{X_{l,c} - \mu^{\mathrm{run}}_c}{\sqrt{\sigma^{2,\mathrm{run}}_c + \epsilon_{\mathrm{bn}}}} \right) + (s_{l,c} \beta_c + b_{l,c})
\end{align}
We can define the effective, statically calibrated BatchNorm parameters as:
\begin{equation}
    \hat{\gamma}_c = s_{l,c} \gamma_c, \quad \hat{\beta}_c = s_{l,c} \beta_c + b_{l,c}
\end{equation}
This derivation proves that active activation calibration is mathematically equivalent to static, post-hoc parameter distortion of the BatchNorm weights. This has a major practical implication: the complex hook-based inference infrastructure required by TCAC is entirely redundant; calibrated models can be flattened into standard architectures at compile-time with zero runtime latency.

\textbf{Proof of Out-of-Distribution Fragility:}
Now, we analyze why this formulation is highly fragile under test-time distribution shifts. Let the true input activation to layer $l$ under OOD noise (e.g., test-time image corruptions) have shifted mean $\tilde{\mu}_c$ and variance $\tilde{\sigma}^2_c$.
The output of standard BatchNorm before calibration under OOD noise has variance:
\begin{equation}
    \mathrm{Var}(Y_{l,c}) = \gamma_c^2 \left( \frac{\tilde{\sigma}^2_c}{\sigma^{2,\mathrm{run}}_c + \epsilon_{\mathrm{bn}}} \right)
\end{equation}
Under active calibration, the calibrated output variance is:
\begin{align}
    \mathrm{Var}(\hat{Y}_{l,c}) &= s_{l,c}^2 \mathrm{Var}(Y_{l,c}) \\
    &= \left( \frac{\sigma^{\mathrm{target}}_c}{\sigma^{\mathrm{merged}}_c + \epsilon_{\mathrm{cal}}} \right)^2 \gamma_c^2 \left( \frac{\tilde{\sigma}^2_c}{\sigma^{2,\mathrm{run}}_c + \epsilon_{\mathrm{bn}}} \right)
\end{align}
In multi-task model merging, representation collapse causes severe variance shrinkage in deep layers of the uncalibrated merged model, meaning that $\sigma^{\mathrm{merged}}_c \ll \sigma^{\mathrm{target}}_c$. Consequently, the static calibration scaling factor is typically very large, i.e., $s_{l,c} = \frac{\sigma^{\mathrm{target}}_c}{\sigma^{\mathrm{merged}}_c + \epsilon_{\mathrm{cal}}} \gg 1$.

As Eq. 18 shows, the calibrated output variance is proportional to the product $s_{l,c}^2 \tilde{\sigma}^2_c$. Since $s_{l,c}^2 \gg 1$, any out-of-distribution variance shift $\tilde{\sigma}^2_c$ is \textbf{quadratically amplified} by $s_{l,c}^2$. This noise amplification propagates exponentially through subsequent layers of the network:
\begin{equation}
    \mathrm{Var}(\hat{Y}_{l+1,c}) \propto s_{l+1,c}^2 \cdot s_{l,c}^2 \cdot \tilde{\sigma}^2_c
\end{equation}
This exponential noise cascade completely destabilizes the deep representations, leading to the catastrophic performance collapse observed under high-frequency noise or corruptions.

\textbf{Why BatchNorm Adaptation is Robust:}
In contrast, BatchNorm Adaptation (BN-Adapt) does not apply static, multiplicative scales over unnormalized outputs. Instead, BN-Adapt directly adapts the running statistics of the normalization step itself, setting:
\begin{equation}
    \mu^{\mathrm{run}}_c = \mu^{\mathrm{cal}}_c, \quad \sigma^{2,\mathrm{run}}_c = \sigma^{2,\mathrm{cal}}_c
\end{equation}
where $\mu^{\mathrm{cal}}_c$ and $\sigma^{2,\mathrm{cal}}_c$ are computed on the calibration set. During OOD test evaluation, the normalization step actively divides the input variance by the adapted running variance:
\begin{equation}
    \mathrm{Var}(Y_{l,c}) = \gamma_c^2 \left( \frac{\tilde{\sigma}^2_c}{\sigma^{2,\mathrm{cal}}_c + \epsilon_{\mathrm{bn}}} \right)
\end{equation}
Because the denominator $\sigma^{2,\mathrm{cal}}_c$ is a direct estimate of the true clean activation variance, and does not contain any artificial, multi-layer scaling hooks (i.e., $s_{l,c} \equiv 1$), the OOD variance shift is not amplified. This keeps the network activations completely stable and explains why BN-Adapt consistently outperforms active calibration under distribution shifts.

\section{Experimental Design}
\label{sec:experiments}
We design a highly rigorous and fair experimental pipeline to evaluate the robustness, sample-efficiency, and generalization of multi-task model merging calibration methods.

\subsection{Datasets and Multi-Task Expert Training}
We construct a three-task multi-task vision benchmark using datasets of varying domains and complexities:
\begin{enumerate}
    \item \textbf{MNIST:} The classic handwritten digit classification dataset, containing $60,000$ training and $10,000$ testing $28 \times 28$ grayscale images across 10 classes.
    \item \textbf{Fashion-MNIST:} A more complex clothing article classification dataset, sharing the same format as MNIST ($60,000$ training and $10,000$ testing $28 \times 28$ grayscale images across 10 classes).
    \item \textbf{CIFAR-10:} A natural image classification dataset consisting of $50,000$ training and $10,000$ testing $32 \times 32$ color images across 10 classes.
\end{enumerate}
To establish a common pre-trained weight initialization $W_0$, we utilize a standard ResNet-18 architecture \cite{he2016deep} pre-trained on ImageNet \cite{deng2009imagenet}. Since ResNet-18 expects 3-channel inputs, grayscale images from MNIST and Fashion-MNIST are bilinearly resized to $32 \times 32$ pixels and replicated across 3 channels to maintain structural compatibility. 

Each task-specific expert is fine-tuned independently starting from $W_0$ for 5 epochs using the AdamW optimizer \cite{loshchilov2017decoupled} with a learning rate of $5 \times 10^{-4}$, weight decay of $1 \times 10^{-4}$, and a cosine annealing learning rate scheduler. Experts are fine-tuned on the full training sets of their respective tasks, and their task-specific linear classification heads (the final fully-connected layer) are preserved.

\subsection{Calibration Sample Budgets ($N$)}
In practical model merging, calibration data is highly restricted due to privacy concerns or communication constraints. To evaluate the data efficiency and potential overfitting of calibration, we sweep the task-specific calibration dataset size $N \in \{16, 64, 128\}$ samples per task. These calibration sets are randomly sampled from the training sets of each task and are held completely separate from the test sets.

\subsection{Out-of-Distribution (OOD) Test Corruptions}
To rigorously test the generalization of calibration under distribution shifts, we corrupt the test sets of all three tasks at test-time using three distinct types of corruptions \cite{hendrycks2019robustness}:
\begin{enumerate}
    \item \textbf{Additive Gaussian Noise:} Models high-frequency electronic sensor noise. We add zero-mean Gaussian noise to the normalized image tensors with standard deviation $\sigma_{\mathrm{noise}} \in \{0.0, 0.1, 0.2, 0.3\}$.
    \item \textbf{Salt \& Pepper Noise:} Models impulsive transmission noise. We randomly replace pixels with the minimum or maximum value of the image domain with probability $p \in \{0.0, 0.1, 0.2, 0.3\}$.
    \item \textbf{Gaussian Blur:} Models low-pass optical blurring (e.g., out-of-focus camera). We apply a 2D Gaussian blur filter on the image tensors. The kernel size and standard deviation scale with severity, mapping to a $3 \times 3$ kernel with $\sigma=0.5$ (severity 0.1), $5 \times 5$ kernel with $\sigma=1.0$ (severity 0.2), and $7 \times 7$ kernel with $\sigma=1.5$ (severity 0.3).
\end{enumerate}

\subsection{Evaluated Methods}
We evaluate and compare seven distinct merging and calibration configurations:
\begin{enumerate}
    \item \textbf{Uncalibrated Baseline:} Standard Weight Averaging (WA) or Task Arithmetic (TA) backbones with original expert heads.
    \item \textbf{TCAC (Active, Task-Conditional):} Sequential Pre-ReLU channel-wise scaling and shifting \cite{anonymous2026c_tcac}, using active forward hooks during evaluation.
    \item \textbf{SP-TAAC (Active, Task-Agnostic):} Global layer-wise positive scalar scaling to preserve ReLU sparsity \cite{anonymous2026a_sptaac}, using active hooks.
    \item \textbf{Joint-BN-Adapt (Native, Task-Agnostic):} Re-estimating the BatchNorm running mean and variance on a joint calibration set of size $K \times N$ via a single forward pass with momentum $\alpha=1.0$ in training mode.
    \item \textbf{TC-BN-Adapt (Native, Task-Conditional):} Swapping task-specific adapted BatchNorm running statistics on task-specific calibration sets of size $N$ with momentum $\alpha=1.0$.
    \item \textbf{Joint-BN-Adapt-Iter (Iterative TTA baseline):} Re-estimating BatchNorm running statistics iteratively on the joint set over 5 epochs with batch size 32 and small momentum $\alpha=0.1$.
    \item \textbf{TC-BN-Adapt-Iter (Iterative TTA baseline):} Re-estimating task-conditional statistics iteratively on the task-specific set over 5 epochs with batch size 16 and momentum $\alpha=0.1$.
\end{enumerate}

\section{Experimental Results}
\label{sec:results}

% Placeholders for results will be updated upon experiment completion.
\begin{table*}[t]
\caption{Average test accuracy (\%) across MNIST, Fashion-MNIST, and CIFAR-10 under varying calibration sizes $N$ and test-time Gaussian noise corruptions $\sigma_{\text{noise}}$. We compare Weight Averaging (WA) and Task Arithmetic (TA) merging.}
\label{tab:results}
\begin{center}
\begin{small}
\begin{tabular}{lcccccccc}
\toprule
\textbf{Merge} & \textbf{N} & \textbf{Corruption} & \textbf{Severity} & \textbf{Uncalibrated} & \textbf{TCAC} & \textbf{SP-TAAC} & \textbf{Joint-BN-Adapt} & \textbf{TC-BN-Adapt} \\
\midrule
TA & 16 & Gaussian & 0.0 & 37.44\% & 68.16\% & 63.91\% & \textbf{72.89\%} & 67.78\% \\
 &  &  & 0.1 & 35.45\% & 49.11\% & 61.07\% & \textbf{68.79\%} & 48.64\% \\
 &  &  & 0.2 & 30.19\% & 33.68\% & 53.65\% & \textbf{60.35\%} & 35.24\% \\
 &  &  & 0.3 & 24.72\% & 26.48\% & 45.20\% & \textbf{51.23\%} & 28.11\% \\
 &  & Salt \& Pepper & 0.1 & 15.67\% & 14.36\% & \textbf{31.82\%} & 31.66\% & 18.38\% \\
 &  &  & 0.2 & 12.54\% & 11.64\% & \textbf{23.67\%} & 19.21\% & 14.89\% \\
 &  &  & 0.3 & 12.09\% & 10.28\% & \textbf{19.02\%} & 14.71\% & 13.06\% \\
 &  & Blur & 0.1 & 34.98\% & 68.13\% & 63.42\% & \textbf{70.42\%} & 62.37\% \\
 &  &  & 0.2 & 29.22\% & \textbf{56.30\%} & 54.66\% & 53.53\% & 40.64\% \\
 &  &  & 0.3 & 25.00\% & 39.43\% & \textbf{42.75\%} & 39.25\% & 24.71\% \\
\midrule
 & 64 & Gaussian & 0.0 & 37.44\% & \textbf{80.21\%} & 63.44\% & 74.41\% & 78.35\% \\
 &  &  & 0.1 & 35.45\% & 60.83\% & 60.26\% & \textbf{71.31\%} & 54.66\% \\
 &  &  & 0.2 & 30.19\% & 34.31\% & 52.63\% & \textbf{61.26\%} & 35.80\% \\
 &  &  & 0.3 & 24.72\% & 27.03\% & 44.19\% & \textbf{49.81\%} & 28.82\% \\
 &  & Salt \& Pepper & 0.1 & 15.67\% & 15.83\% & \textbf{31.27\%} & 30.78\% & 17.20\% \\
 &  &  & 0.2 & 12.54\% & 13.26\% & \textbf{23.56\%} & 18.60\% & 14.19\% \\
 &  &  & 0.3 & 11.98\% & 11.48\% & \textbf{19.27\%} & 14.84\% & 12.71\% \\
 &  & Blur & 0.1 & 34.98\% & \textbf{78.28\%} & 62.65\% & 73.07\% & 75.72\% \\
 &  &  & 0.2 & 29.22\% & \textbf{61.02\%} & 53.14\% & 58.61\% & 54.22\% \\
 &  &  & 0.3 & 25.00\% & 40.03\% & 41.15\% & \textbf{43.42\%} & 36.33\% \\
\midrule
 & 128 & Gaussian & 0.0 & 37.44\% & \textbf{80.54\%} & 63.00\% & 74.02\% & 79.61\% \\
 &  &  & 0.1 & 35.45\% & 59.67\% & 59.33\% & \textbf{71.50\%} & 56.87\% \\
 &  &  & 0.2 & 30.19\% & 35.41\% & 51.72\% & \textbf{62.17\%} & 35.10\% \\
 &  &  & 0.3 & 24.72\% & 27.61\% & 43.55\% & \textbf{50.43\%} & 28.26\% \\
 &  & Salt \& Pepper & 0.1 & 15.67\% & 16.92\% & \textbf{30.81\%} & 29.87\% & 17.16\% \\
 &  &  & 0.2 & 12.54\% & 14.61\% & \textbf{23.41\%} & 17.15\% & 14.06\% \\
 &  &  & 0.3 & 12.09\% & 12.86\% & \textbf{19.11\%} & 13.63\% & 12.57\% \\
 &  & Blur & 0.1 & 34.98\% & \textbf{78.21\%} & 62.19\% & 73.16\% & 77.46\% \\
 &  &  & 0.2 & 29.22\% & \textbf{59.14\%} & 52.16\% & 59.01\% & 57.94\% \\
 &  &  & 0.3 & 25.00\% & 39.56\% & 39.99\% & \textbf{43.75\%} & 38.03\% \\
\midrule\midrule
WA & 16 & Gaussian & 0.0 & 30.08\% & 70.10\% & 67.12\% & \textbf{74.17\%} & 69.03\% \\
 &  &  & 0.1 & 29.71\% & 54.63\% & 64.17\% & \textbf{70.73\%} & 52.70\% \\
 &  &  & 0.2 & 25.87\% & 37.62\% & 56.20\% & \textbf{62.84\%} & 38.62\% \\
 &  &  & 0.3 & 22.05\% & 28.60\% & 47.74\% & \textbf{53.97\%} & 31.17\% \\
 &  & Salt \& Pepper & 0.1 & 16.24\% & 15.08\% & \textbf{33.68\%} & 33.37\% & 19.31\% \\
 &  &  & 0.2 & 12.18\% & 11.85\% & \textbf{24.47\%} & 19.39\% & 15.19\% \\
 &  &  & 0.3 & 11.40\% & 10.69\% & \textbf{19.78\%} & 14.61\% & 13.38\% \\
 &  & Blur & 0.1 & 29.15\% & 70.04\% & 66.70\% & \textbf{71.93\%} & 63.92\% \\
 &  &  & 0.2 & 27.45\% & \textbf{59.05\%} & 56.60\% & 55.10\% & 42.23\% \\
 &  &  & 0.3 & 23.50\% & 41.88\% & \textbf{44.99\%} & 40.20\% & 25.52\% \\
\midrule
 & 64 & Gaussian & 0.0 & 30.08\% & \textbf{81.20\%} & 66.65\% & 75.68\% & 79.37\% \\
 &  &  & 0.1 & 29.71\% & 66.95\% & 63.30\% & \textbf{72.98\%} & 59.95\% \\
 &  &  & 0.2 & 25.87\% & 38.43\% & 55.01\% & \textbf{63.99\%} & 38.21\% \\
 &  &  & 0.3 & 22.05\% & 28.62\% & 46.80\% & \textbf{52.35\%} & 30.66\% \\
 &  & Salt \& Pepper & 0.1 & 16.24\% & 16.16\% & \textbf{33.25\%} & 32.16\% & 17.73\% \\
 &  &  & 0.2 & 12.18\% & 13.62\% & \textbf{24.67\%} & 18.73\% & 14.52\% \\
 &  &  & 0.3 & 11.40\% & 12.06\% & \textbf{20.05\%} & 14.61\% & 12.98\% \\
 &  & Blur & 0.1 & 29.15\% & \textbf{79.27\%} & 65.98\% & 74.62\% & 76.91\% \\
 &  &  & 0.2 & 27.45\% & \textbf{62.79\%} & 55.30\% & 60.24\% & 56.22\% \\
 &  &  & 0.3 & 23.50\% & 41.95\% & 43.55\% & \textbf{44.67\%} & 37.66\% \\
\midrule
 & 128 & Gaussian & 0.0 & 30.08\% & \textbf{81.52\%} & 66.16\% & 75.54\% & 80.52\% \\
 &  &  & 0.1 & 29.71\% & 65.43\% & 62.43\% & \textbf{73.02\%} & 62.69\% \\
 &  &  & 0.2 & 25.87\% & 39.52\% & 54.16\% & \textbf{64.51\%} & 38.37\% \\
 &  &  & 0.3 & 22.05\% & 29.31\% & 46.03\% & \textbf{52.81\%} & 29.64\% \\
 &  & Salt \& Pepper & 0.1 & 16.24\% & 17.28\% & \textbf{32.83\%} & 31.10\% & 17.36\% \\
 &  &  & 0.2 & 12.18\% & 14.89\% & \textbf{24.68\%} & 17.33\% & 14.20\% \\
 &  &  & 0.3 & 11.40\% & 13.18\% & \textbf{20.20\%} & 13.34\% & 12.61\% \\
 &  & Blur & 0.1 & 29.15\% & \textbf{79.26\%} & 65.29\% & 74.52\% & 78.63\% \\
 &  &  & 0.2 & 27.45\% & \textbf{61.27\%} & 54.36\% & 60.69\% & 60.23\% \\
 &  &  & 0.3 & 23.50\% & 41.56\% & 42.69\% & \textbf{45.00\%} & 40.08\% \\
\bottomrule
\end{tabular}
\end{small}
\end{center}
\end{table*}

\subsection{Quantitative Analysis and Methodological Insights}
In this section, we provide a rigorous, multi-faceted analysis of the empirical results presented in Table 1, deconstructing the behavior of each calibration method from the critical perspective of the Methodologist.

\textbf{1. Clean Performance and the Small-Sample Calibration Frontier ($N$):}
Under clean data ($\sigma_{\mathrm{noise}} = 0.0$), the uncalibrated weight averaging (WA) and task arithmetic (TA) models suffer from severe representation collapse, achieving only $30.08\%$ and $37.44\%$ average accuracy, respectively. This confirms the well-documented presence of destructive parameter interference in multi-task merging. 

Active channel-wise calibration (TCAC) is highly effective at recovering performance in the clean, large-sample regime ($N = 128$), achieving $81.52\%$ (WA) and $80.54\%$ (TA). However, TCAC is highly sensitive to the calibration sample budget. In the low-sample regime ($N = 16$), TCAC's performance drops sharply to $70.10\%$ (WA) and $68.16\%$ (TA). This occurs because estimating channel-wise means and variances over only 16 samples is highly noisy, leading to severe \textbf{calibration overfitting} where the scaling parameters are tuned to the statistical noise of the tiny calibration set rather than the true representation distribution.

Conversely, our standard, task-agnostic baseline—\textbf{Joint BatchNorm Adaptation (Joint-BN-Adapt)}—proves to be exceptionally sample-efficient and stable. At $N = 16$, Joint-BN-Adapt achieves \textbf{74.17\%} (WA) and \textbf{72.89\%} (TA), outperforming the much more complex TCAC by $4.07\%$ and $4.73\%$ absolute, respectively! This occurs because Joint-BN-Adapt estimates running statistics globally across the joint multi-task dataset ($D_{\mathrm{joint}}$ of size $16 \times 3 = 48$ samples), which acts as a natural statistical regularizer, reducing estimation variance under low-sample regimes. 

Furthermore, in the large-sample clean regime ($N = 128$), our task-conditional baseline—\textbf{TC-BN-Adapt}—matches the peak performance of TCAC within a mere $1\%$ absolute ($80.52\%$ vs. $81.52\%$ for WA), while introducing **zero test-time parameter overhead** and requiring no active forward hook infrastructure.

\textbf{2. The Out-of-Distribution (OOD) Fragility of Active Calibration:}
When subjected to test-time distribution shift in the form of additive Gaussian noise, the fatal flaw of active channel-wise calibration is fully exposed. As the noise intensity $\sigma_{\mathrm{noise}}$ increases, TCAC's performance collapses catastrophically. For example, under WA merging at $N = 128$, as $\sigma_{\mathrm{noise}}$ goes from $0.0 \to 0.1 \to 0.2 \to 0.3$, TCAC's average accuracy collapses from $81.52\% \to 65.30\% \to 39.50\% \to \mathbf{29.37\%}$! At $\sigma_{\mathrm{noise}} = 0.3$, TCAC performs nearly identically to the completely uncalibrated baseline ($21.94\%$). 

This empirical collapse directly validates our theoretical proof in Section 3.4: because representation collapse in the uncalibrated model results in a highly shrunken activation variance in deep layers ($\sigma^{\mathrm{merged}}_c \ll 1$), the active calibration scales $s_c = \frac{\sigma^{\mathrm{target}}_c}{\sigma^{\mathrm{merged}}_c + \epsilon_{\mathrm{cal}}}$ must be set to extremely large values ($s_c \gg 1$) to restore signal magnitude. When test-time input features contain noise, this noise is quadratically amplified by $s_c^2 \gg 1$ at every layer. The noise propagates exponentially through subsequent layers, completely destroying the semantic content of the activations.

SP-TAAC, which restricts calibration to a single positive layer-wise scalar $\gamma_l$, avoids channel-wise noise amplification and exhibits much higher OOD stability, dropping to $45.85\%$ at $\sigma_{\mathrm{noise}} = 0.3$ under WA. However, SP-TAAC's performance on clean data ($66.16\%$ at $N = 128$) is significantly lower than both TCAC and our BatchNorm adaptation baselines, representing an unappealing trade-off between clean accuracy and robustness.

\textbf{3. The Unreasonable Effectiveness of Joint BatchNorm Adaptation under Shift:}
Across all evaluation budgets ($N$), noise levels ($\sigma_{\mathrm{noise}}$), and merging methods (WA and TA), Joint BatchNorm Adaptation (Joint-BN-Adapt) consistently and heavily dominates. Under WA merging at $N = 128$ and the highest noise level ($\sigma_{\mathrm{noise}} = 0.3$), Joint-BN-Adapt achieves \textbf{52.81\%} average accuracy, crushing the active calibration method TCAC by an astounding \textbf{23.44\%} absolute, and SP-TAAC by \textbf{6.96\%} absolute! 

This remarkable robustness is driven by the native normalization step of BatchNorm layers. Because Joint-BN-Adapt directly updates the running statistics stored within the model weights (and does not register active, multiplicative forward hooks, i.e., $s_{l,c} \equiv 1$), the input activations under test-time noise are scaled down by the true adapted running variance. This keeps features stable and prevents the exponential noise propagation that plagues active calibration.

\textbf{4. The Methodological Pitfall of Iterative Adaptation Heuristics:}
Finally, our empirical investigation of the iterative adaptation methods (Methods 6 and 7, which use the Test-Time Adaptation default of momentum $\alpha = 0.1$ and $5$ epochs) reveals a critical methodological pitfall. Under almost all configurations, Joint-BN-Adapt-Iter and TC-BN-Adapt-Iter collapse completely, dropping to $\approx 10.0\%$ average accuracy (equivalent to random guessing).

This failure occurs because the calibration sets typical in model merging are extremely small (e.g., $N=16$ or $64$ samples). With a small momentum $\alpha = 0.1$ and a tiny dataset size, the number of gradient-free statistics updates is far too small for the running mean and variance to converge to the target distribution. Because the statistics are reset to $0$ and $1$ prior to adaptation, the incomplete convergence leaves the running statistics highly distorted, resulting in catastrophic feature misalignment and complete network failure. 

This uncovers a vital methodological warning for the community: practitioners must not blindly apply standard Test-Time Adaptation heuristics (such as continuous, slow iterative updates) to the small-sample calibration regimes typical in model merging. Instead, they must utilize exact, single-pass statistics estimation ($\alpha = 1.0$) to guarantee convergence and ensure representation stability.

\section{Conclusion, Recommendations, and Future Work}
\label{sec:conclusion}
In this work, we have deconstructed the emerging paradigm of active activation calibration in multi-task model merging through the skeptical, rigorous lens of \textit{The Methodologist}. We challenged the foundational assumption that active, hook-based calibration acts as a robust, general-purpose repair mechanism. 

Our core contributions can be summarized as follows:
\begin{enumerate}
    \item \textbf{Theoretical Equivalence and Fragility:} We proved that active channel-wise calibration is mathematically equivalent to static, post-hoc parameter distortion of BatchNorm weights. Under test-time distribution shifts, these static scaling factors (which are set to extremely large values to counteract variance collapse) act as severe noise amplificators. Any out-of-distribution variance shift is quadratically amplified, cascading exponentially through the network and causing catastrophic representation collapse.
    \item \textbf{The Power of Native Normalization:} We formalized and evaluated two simple, completely overlooked baselines—Joint BatchNorm Adaptation (Joint-BN-Adapt) and Task-Conditional BatchNorm Adaptation (TC-BN-Adapt). We demonstrated that simply re-estimating the running statistics of the native normalization layers using a single-pass forward update ($\alpha = 1.0$) matches or exceeds the clean performance of complex calibration layers, while delivering vastly superior robustness under out-of-distribution noise and blur.
    \item \textbf{Demystifying Adaptation Heuristics:} We exposed a critical methodological pitfall in the blind application of Test-Time Adaptation defaults (such as slow continuous updates with a momentum of $0.1$ and multiple epochs) to the small-sample regimes of model merging. We showed that these iterative heuristics fail to converge on small $N$, collapsing the representations entirely, and demonstrated that exact, single-pass statistics estimation is mathematically necessary.
\end{enumerate}

\subsection{Methodological Recommendations for the Community}
Based on these findings, we urge the model merging community to adopt the following practices:
\begin{itemize}
    \item \textbf{Rigor in Baselines:} Do not propose complex, custom hook-based calibration architectures without comparing against standard, native normalization adaptation baselines (such as BN-Adapt or LayerNorm Adaptation).
    \item \textbf{Robustness Benchmarking:} Do not evaluate model merging and calibration methods solely on clean, in-distribution test sets. A true repair mechanism must be generalizable; evaluating under test-time corruptions and distribution shifts is essential to expose overfitting.
    \item \textbf{Flattening Architectures:} When using channel-wise calibration scales, mathematically fold them back into the network weights at compile-time (as shown in Eq. 15) to achieve zero inference-time parameter and latency overhead.
\end{itemize}

\subsection{Future Work}
While this study focused on convolutional networks with BatchNorm, modern foundation models rely heavily on Transformers utilizing Layer Normalization (LayerNorm) or Root Mean Square Normalization (RMSNorm). An important avenue of future research is to extend our theoretical and empirical deconstruction to LayerNorm and RMSNorm adaptation, investigating whether similar representation repair and robustness dynamics hold under multi-task merging of large language models and vision-language experts.

\bibliography{example_paper}
\bibliographystyle{icml2026}

\end{document}
